\section{Referencial teórico}

Esta seção organiza os fundamentos teóricos que sustentam a pergunta de pesquisa e a delimitação conceitual do projeto. O argumento está estruturado em três eixos complementares: a relação entre inteligência artificial generativa, RAG e recuperação da informação; as especificidades da informação legislativa como acervo público; e os critérios de avaliação aplicáveis a sistemas RAG em domínios institucionais.

\subsection{Inteligência artificial generativa, RAG e recuperação da informação}

Os grandes modelos de linguagem ampliaram a capacidade de interação com sistemas computacionais por meio de linguagem natural, viabilizando tarefas de sumarização, classificação, geração de texto, resposta a perguntas e apoio à análise documental \parencite{brownEtAl2020fewShotLearners,zhaoEtAl2023llmSurvey,minaee2024llmSurvey}. A arquitetura dominante desses modelos se apoia em avanços associados aos \textit{transformers}, combinados a estratégias de pré-treinamento em larga escala e ajuste para seguir instruções \parencite{vaswani2017attention,brownEtAl2020fewShotLearners,ouyangEtAl2022instructGpt}. Apesar dessas capacidades, LLMs não são mecanismos de recuperação documental por definição: suas respostas podem refletir padrões aprendidos no treinamento, inferências probabilísticas e lacunas de atualização, sem indicar de forma confiável quais documentos sustentam cada afirmação \parencite{jiEtAl2023hallucinationSurvey,anhHoang2025hallucinations,dahlEtAl2024largeLegalFictions}.

O RAG surge como uma estratégia para aproximar geração textual e recuperação da informação. Na formulação clássica, a arquitetura combina um mecanismo de recuperação sobre uma base documental externa com um modelo generativo que produz a resposta a partir do contexto recuperado \parencite{lewisEtAl2020rag}. Revisões recentes classificam essas arquiteturas em configurações ingênuas, avançadas e modulares, destacando que o desempenho depende de todo o pipeline, não apenas do modelo final de geração \parencite{gao2023ragSurvey,gao2024modularRag}. Componentes como extração documental, segmentação em \textit{chunks}, enriquecimento por metadados, geração de \textit{embeddings}, recuperação vetorial, reranqueamento, seleção de contexto e \textit{prompting} orientado a evidências influenciam diretamente a qualidade das respostas \parencite{gupta2024comprehensiveRag,nie2024textEmbeddingSurvey}.

No componente de recuperação, o problema se aproxima da tradição de recuperação da informação, que envolve representação de documentos, formulação de consultas, ranqueamento e avaliação da pertinência dos resultados recuperados \parencite{manningEtAl2008informationRetrieval}. Em sistemas RAG, essa etapa é frequentemente operacionalizada por recuperação densa ou vetorial, em que representações semânticas são usadas para localizar passagens relevantes ao contexto da consulta \parencite{karpukhinEtAl2020dpr,thakurEtAl2021beir}.

Essa literatura é central para o projeto porque desloca a análise de uma discussão genérica sobre ``uso de IA'' para uma questão avaliativa delimitada: em que condições uma arquitetura RAG recupera evidências pertinentes, gera respostas verificáveis e reconhece limites de escopo. Essa delimitação exige distinguir, no desenho da pesquisa, a recuperação documental, a geração de respostas e o julgamento das respostas.

\subsection{Informação legislativa, transparência e acervos parlamentares}

O acesso à informação pública é dimensão central da administração democrática, mas a disponibilidade formal de documentos não elimina barreiras de compreensão, localização e uso qualificado. No domínio legislativo, discursos, debates e registros parlamentares são fontes relevantes para reconstruir posicionamentos, agendas, conflitos, justificativas e memória institucional \parencite{skubicFiser2024discourseReview}. Ao mesmo tempo, esses acervos podem ser difíceis de explorar por mecanismos convencionais de busca, sobretudo quando o usuário não conhece previamente a terminologia exata, os atores envolvidos ou a organização dos registros \parencite{bandeiraBernardes2016information,geunis2023parliamentaryMonitoring}.

No caso brasileiro, estudos recentes da Consultoria Legislativa do Senado mostram que os pronunciamentos em plenário formam um corpus extenso, dinâmico e analiticamente relevante, com variações de volume, extensão, interatividade e complexidade textual ao longo do tempo \parencite{blanco2025plenarioPalanqueEstudio,blanco2026letrasGraudas}. Em paralelo, pesquisas sobre organização do conhecimento legislativo indicam que o uso de inteligência artificial nesse domínio precisa considerar estrutura informacional, explicabilidade, mediação documental e confiabilidade dos resultados \parencite{martelloteViola2025organizacaoConhecimento}.

Assim, a informação legislativa não deve ser tratada apenas como massa textual a ser processada. Trata-se de um recurso institucional, marcado por regras de produção, metadados, autoria, temporalidade, retificações, usos políticos e requisitos de accountability \parencite{bandeiraBernardes2016information,martelloteViola2025organizacaoConhecimento}. Essa característica reforça a necessidade de combinar inovação técnica com critérios de avaliação compatíveis com o uso público e institucional da informação legislativa.

\subsection{Avaliação de sistemas RAG em domínios institucionais}

A avaliação de sistemas RAG apresenta desafios próprios porque uma resposta final aparentemente correta pode ocultar recuperação deficiente, uso inadequado das fontes, extrapolação além do contexto ou julgamento automatizado permissivo. Essa dificuldade se relaciona ao caráter multidimensional da avaliação de sistemas baseados em modelos de linguagem, nos quais diferentes dimensões de desempenho, confiabilidade e uso precisam ser analisadas de forma articulada \parencite{liangEtAl2023helm}. Por isso, a literatura recomenda distinguir métricas de recuperação, critérios de avaliação da geração, avaliação humana, avaliação automatizada baseada em LLMs e análise qualitativa de erros \parencite{saadFalconEtAl2024ares,esEtAl2023ragas,pipitoneAlami2024legalbenchRag}. Benchmarks recentes também apontam a importância de testar diferentes famílias de perguntas, como questões factuais, comparativas, temporais, multietapas, numéricas e não respondíveis \parencite{zhuEtAl2025rageval}.

Entre essas técnicas, o paradigma \textit{LLM-as-a-judge} merece atenção específica neste projeto porque permite avaliar respostas geradas com base em rubricas explícitas \parencite{liuEtAl2023geval,liuEtAl2025judgeRag}. Essa abordagem pode ajudar a escalar avaliações, mas não elimina a necessidade de validação humana, pois o julgamento automatizado pode variar conforme modelo, rubrica, formato de entrada e complexidade da tarefa. Em domínios públicos, esse ponto é decisivo: uma avaliação automatizada não deve ser confundida com certificação de qualidade institucional.

Em domínios públicos, diretrizes e estudos sobre IA enfatizam transparência, finalidade pública, supervisão humana e responsabilização \parencite{ipu2024aiParliaments,wfd2023guidelinesAiParliaments,deAlmeidaSantos2025aiGovernance}. Nesta pesquisa, esses elementos são tratados como contexto institucional para justificar a necessidade de avaliação criteriosa, e não como dimensão autônoma do protocolo. O protocolo concentra-se em aspectos diretamente observáveis no sistema RAG: recuperação documental, geração de respostas, indicação das fontes, explicitação de insuficiência documental e verificabilidade das fontes.

Nesse contexto, protocolo é entendido como uma organização explícita de procedimentos, métricas, rubricas e critérios de interpretação voltados a avaliar, de forma reprodutível, dimensões complementares de um sistema RAG.

\subsection{Lacuna de pesquisa}

A literatura existente oferece bases técnicas para compreender arquiteturas RAG, avaliação de respostas geradas e recuperação da informação, bem como estudos sobre informação legislativa e parlamentos. Entretanto, os trabalhos mais próximos distribuem-se em frentes ainda pouco integradas. Estudos sobre RAG enfatizam arquitetura, recuperação, geração e avaliação automatizada, mas frequentemente tratam domínios genéricos ou jurídicos mais amplos \parencite{gao2023ragSurvey,saadFalconEtAl2024ares,esEtAl2023ragas,pipitoneAlami2024legalbenchRag}. Pesquisas sobre parlamentos e informação legislativa discutem acesso, transparência, discurso parlamentar e organização do conhecimento, mas nem sempre implementam ou avaliam sistemas conversacionais baseados em RAG \parencite{skubicFiser2024discourseReview,blanco2025plenarioPalanqueEstudio,blanco2026letrasGraudas,martelloteViola2025organizacaoConhecimento}. Aplicações recentes em legislação e política, por sua vez, exploram fact-checking, redação legislativa, consulta a textos normativos, interfaces conversacionais e arquiteturas híbridas, mas diferem quanto a corpus, língua, jurisdição, tarefa e desenho de avaliação \parencite{khaliq2024ragar,chouhanGertz2024lexdrafter,garrisonEtAl2024talkNdaa,rogiersEtAl2024kamerraad,matoshiEtAl2025parliamentRag,colomboEtAl2025legisSearch,mosbachEtAl2025worldAvatarParliament}.

Desse modo, a lacuna refere-se à escassez de estudos empíricos que combinem, no mesmo desenho de pesquisa e no contexto brasileiro, um acervo parlamentar real, um \textit{pipeline} RAG reprodutível, avaliação formal da recuperação documental, avaliação das respostas geradas e análise da verificabilidade das fontes. No caso desta pesquisa, essa lacuna será examinada em discursos da 56\textordfeminine{} Legislatura do Senado Federal, com atenção simultânea à recuperação, à geração e à verificabilidade das fontes. O foco do estudo recai sobre a forma como a recuperação documental, a geração de respostas e a verificabilidade das fontes podem ser avaliadas de maneira articulada em consultas legislativas mediadas por RAG.

Essa lacuna é especialmente relevante para a administração pública, pois envolve simultaneamente organização de um protocolo avaliativo, avaliação empírica e uso de acervos institucionais de interesse público. A dissertação proposta busca contribuir para esse espaço ao desenvolver uma investigação com critérios explícitos para avaliação de sistemas RAG em acervos legislativos.
